In practical code, some buffer declarations must specify a \textbf{storage specifier}, like so:

\begin{lstlisting}[language=croqtile,caption={Storage Specifier}]
global f32 [32, 7, 2] a;
shared u8 [512, 144] b;
local u8 [72, 1024] c;
\end{lstlisting}

Since Croqtile handles storage in a heterogenous context, a buffer definition without a storage specifier defaults to the storage type of the host program, i.e., CPU memory. Other storage specifiers are defined by the target.

\begin{itemize}[leftmargin=1.6em]
\item \emph{global}: Refers to the device's global storage.
\item \emph{shared}: Refers to the device block's shared storage
\item \emph{local}: Refers to the thread-private storage.
\end{itemize}

Other targets may have different definitions. Furthermore, buffers with different storage types have limitations in their declarations. You will explore this in later chapters.
